\documentclass[12pt,a4paper]{amsart}

\usepackage{amsmath,amssymb,amsthm}
\usepackage{mathtools}
\usepackage[ngerman]{babel}
\usepackage{hyperref}
\usepackage{geometry}
\geometry{margin=2.5cm}

% -------------------------------------------------------
% Theorem-Umgebungen
% -------------------------------------------------------
\newtheorem{theorem}{Satz}[section]
\newtheorem{lemma}[theorem]{Lemma}
\newtheorem{corollary}[theorem]{Korollar}
\newtheorem{proposition}[theorem]{Proposition}
\theoremstyle{definition}
\newtheorem{definition}[theorem]{Definition}
\newtheorem{remark}[theorem]{Bemerkung}
\newtheorem{example}[theorem]{Beispiel}

% -------------------------------------------------------
% Eigene Befehle
% -------------------------------------------------------
\newcommand{\N}{\mathbb{N}}
\newcommand{\Z}{\mathbb{Z}}
\newcommand{\Q}{\mathbb{Q}}
\newcommand{\R}{\mathbb{R}}
\newcommand{\C}{\mathbb{C}}
\DeclareMathOperator{\ord}{ord}

% -------------------------------------------------------
\begin{document}
% -------------------------------------------------------

\title{Kongruente Zahlen, elliptische Kurven und BSD:\\
       Eine konkrete Anwendung der Birch--Swinnerton-Dyer-Vermutung}

\author{Michael Fuhrmann}
\address{Specialist Mathematics Project, Build 84}
\email{specialist-maths@localhost}
\date{\today}

\begin{abstract}
Eine positive ganze Zahl $n$ hei\ss t \emph{kongruente Zahl}, wenn sie
die Fl\"ache eines rechtwinkligen Dreiecks mit rationalen Seiten ist.
Dieses uralte Problem -- arabische Mathematiker im 10.~Jahrhundert --
ist \"aquivalent zur Existenz eines rationalen Punktes unendlicher Ordnung
auf der \emph{kongruenten Zahlkurve}
$E_n: y^2 = x^3 - n^2 x$.

Wir beweisen die \"Aquivalenz, bestimmen die Torsionsuntergruppe
$E_n(\Q)_{\mathrm{tors}} \cong \Z/2\Z \times \Z/2\Z$
und erkl\"aren Tunnells Satz (1983):
Unter Annahme von BSD ist eine quadratfreie ganze Zahl $n$ genau dann
kongruent, wenn
\[
  \#\{(x,y,z) \in \Z^3 : 2x^2 + y^2 + 8z^2 = n\}
  = 2\,\#\{(x,y,z) \in \Z^3 : 2x^2 + y^2 + 32z^2 = n\}
\]
(f\"ur ungerades $n$; eine analoge Formel gilt f\"ur gerades $n$).
Wir verifizieren, dass $5, 6, 7$ kongruent sind,
und diskutieren den Status von $1, 2, 3$.
\end{abstract}

\maketitle
\tableofcontents

% ================================================================
\section{Einleitung: Was ist eine kongruente Zahl?}
% ================================================================

\begin{definition}[Kongruente Zahl]
\label{def:congruent}
Eine positive ganze Zahl $n$ hei\ss t \emph{kongruente Zahl}, wenn es ein
rechtwinkliges Dreieck mit rationalen Seitenl\"angen $a, b, c$
(wobei $c$ die Hypotenuse ist) gibt, dessen Fl\"ache $n$ betr\"agt:
\[
  a^2 + b^2 = c^2, \quad \frac{ab}{2} = n, \quad a, b, c \in \Q_{>0}.
\]
\end{definition}

\begin{example}[Erste Beispiele]
\begin{itemize}
  \item $n = 6$: Das rechtwinklige Dreieck mit Katheten $3, 4$ und
    Hypotenuse $5$ hat Fl\"ache $\tfrac{1}{2} \cdot 3 \cdot 4 = 6$.
    Also ist $6$ eine kongruente Zahl.
  \item $n = 5$: Das Dreieck mit Katheten $\tfrac{3}{2}$, $\tfrac{20}{3}$
    und Hypotenuse $\tfrac{41}{6}$ hat Fl\"ache $5$.
    Also ist $5$ kongruent.
  \item $n = 7$: Das Dreieck mit Katheten $\tfrac{35}{12}$, $\tfrac{24}{5}$
    und Hypotenuse $\tfrac{337}{60}$ hat Fl\"ache $7$.
    Also ist $7$ kongruent.
  \item $n = 1, 2, 3$: Vermutlich nicht kongruent, aber bisher nicht
    bedingungslos bewiesen.
\end{itemize}
\end{example}

\begin{remark}[Historische Notiz]
Kongruente Zahlen wurden von arabischen Mathematikern im 10.~Jahrhundert
studiert und auch von Fibonacci (1225).
Fermat bewies, dass $1$ nicht kongruent ist, mittels unendlichem Abstieg --
der erste dokumentierte Einsatz dieser Methode.
\end{remark}

% ================================================================
\section{Die kongruente Zahlkurve}
% ================================================================

\begin{theorem}[\"Aquivalenz mit elliptischen Kurven]
\label{thm:equiv}
Eine positive quadratfreie ganze Zahl $n$ ist genau dann kongruent,
wenn die elliptische Kurve
\[
  E_n: \quad y^2 = x^3 - n^2 x
\]
einen rationalen Punkt unendlicher Ordnung hat, d.\,h.\
$\mathrm{rang}(E_n(\Q)) \ge 1$.
\end{theorem}

\begin{proof}
\textbf{(Kongruent $\Rightarrow$ Punkt unendlicher Ordnung.)}
Seien $a^2 + b^2 = c^2$ und $ab/2 = n$ mit $a,b,c \in \Q_{>0}$.
Die Standardparametrisierung (vgl.\ \cite{Koblitz1984}, Kapitel~1) liefert
aus $(a,b,c)$ mit $a^2+b^2=c^2$ und $ab=2n$ den Punkt
\[
  P = \Bigl(\bigl(\tfrac{c}{2}\bigr)^2,\; \tfrac{c(a^2-b^2)}{8}\Bigr)
\]
liegt auf $E_n$, denn
\[
  \Bigl(\tfrac{c(a^2-b^2)}{8}\Bigr)^2
  = \frac{c^2(a^2-b^2)^2}{64}
\]
und
\[
  \Bigl(\tfrac{c^2}{4}\Bigr)^3 - n^2\cdot\tfrac{c^2}{4}
  = \frac{c^6}{64} - \frac{n^2 c^2}{4}
  = \frac{c^2}{64}\bigl(c^4 - 16n^2\bigr).
\]
Mit $c^4 - 16n^2 = (a^2+b^2)^2 - (ab)^2 = (a^2-b^2)^2$ folgt
$P \in E_n(\Q)$.

Dieser Punkt liegt nicht in der Torsionsuntergruppe $\{O, (0,0), (\pm n, 0)\}$
(da $y_0 = c(a^2-b^2)/8 \ne 0$ genau dann, wenn $a \ne b$,
d.\,h.\ das rechtwinklige Dreieck nicht gleichschenklig ist), also hat er unendliche Ordnung.

\textbf{(Punkt unendlicher Ordnung $\Rightarrow$ kongruent.)}
Sei $(x_0,y_0) \in E_n(\Q)$ mit $y_0 \ne 0$. Die Zahlen
\[
  a = \left|\frac{x_0^2 - n^2}{y_0}\right|, \quad
  b = \left|\frac{2nx_0}{y_0}\right|, \quad
  c = \left|\frac{x_0^2 + n^2}{y_0}\right|
\]
erf\"ullen $a^2 + b^2 = c^2$ und $ab/2 = n$, wie man mit
$y_0^2 = x_0^3 - n^2 x_0 = x_0(x_0^2-n^2)$ nachrechnet.
\end{proof}

% ================================================================
\section{Die Torsionsuntergruppe von \texorpdfstring{$E_n$}{En}}
% ================================================================

\begin{theorem}[Torsion von $E_n$]
\label{thm:torsion_En}
F\"ur quadratfreies $n \ge 1$:
\[
  E_n(\Q)_{\mathrm{tors}} \;=\; \{\mathcal{O},\, (0,0),\, (n,0),\, (-n,0)\}
  \;\cong\; \Z/2\Z \times \Z/2\Z.
\]
\end{theorem}

\begin{proof}
Die Kurve $E_n: y^2 = x(x-n)(x+n)$ hat drei rationale Wurzeln $0, n, -n$,
die drei $2$-Torsionspunkte $(0,0)$, $(n,0)$, $(-n,0)$ ergeben.
Zusammen mit $\mathcal{O}$ bilden sie $\Z/2\Z \times \Z/2\Z$.

Keine weiteren Torsionspunkte existieren: Nach Nagell--Lutz m\"usste ein
Torsionspunkt $(x_0, y_0)$ mit $y_0 \ne 0$ die Bedingung $y_0^2 \mid \Delta$
(mit $\Delta = 64n^6$) erf\"ullen und $x_0, y_0 \in \Z$ (Nagell--Lutz).
Eine vollst\"andige Fallunterscheidung zeigt, dass keine weiteren
ganzzahligen Punkte auf $E_n$ existieren; vgl.\ \cite{Silverman1992}.
\end{proof}

\begin{corollary}
$n$ ist kongruent $\iff$ $\mathrm{rang}(E_n(\Q)) \ge 1$
$\iff$ $E_n(\Q) \ne E_n(\Q)_{\mathrm{tors}}$.
\end{corollary}

% ================================================================
\section{Tunnells Satz}
% ================================================================

\begin{theorem}[Tunnell, 1983]
\label{thm:tunnell}
F\"ur eine quadratfreie positive ganze Zahl $n$ definiere:
\begin{align*}
  A(n) &= \#\{(x,y,z) \in \Z^3 : 2x^2 + y^2 + 8z^2 = n\}, \\
  B(n) &= \#\{(x,y,z) \in \Z^3 : 2x^2 + y^2 + 32z^2 = n\}, \\
  C(n) &= \#\{(x,y,z) \in \Z^3 : 4x^2 + y^2 + 8z^2 = n/2\}
  \quad (n \text{ gerade}), \\
  D(n) &= \#\{(x,y,z) \in \Z^3 : 4x^2 + y^2 + 32z^2 = n/2\}
  \quad (n \text{ gerade}).
\end{align*}
\begin{enumerate}
  \item[\textup{(i)}] Ist $n$ kongruent, so gilt:
    $A(n) = 2B(n)$ (f\"ur ungerades $n$), bzw.\
    $C(n) = 2D(n)$ (f\"ur gerades $n$).
  \item[\textup{(ii)}] Umgekehrt (unter Annahme von BSD f\"ur $E_n$):
    Falls $A(n) = 2B(n)$ (ungerades $n$) oder $C(n) = 2D(n)$ (gerades $n$),
    dann ist $n$ kongruent.
\end{enumerate}
\end{theorem}

\begin{proof}[Beweisskizze]
Tunnell \cite{Tunnell1983} verwendete Modulformen vom Gewicht $3/2$ und
die Theorie der Theta-Funktionen.
Schritte:
\begin{enumerate}
  \item $L(E_n, 1)$ l\"asst sich durch Werte von Modulformen
    ausdr\"ucken, die tern\"aren quadratischen Formen zugeordnet sind.
  \item Via Shimuras Hebelungsatz und der Theorie halbganzzahliger
    Modulformen ist $L(E_n,1)$ mit der Differenz $A(n)-2B(n)$ verkn\"upft.
  \item Ist $n$ kongruent, so sagt BSD $L(E_n,1) = 0$ vorher,
    was $A(n) = 2B(n)$ ergibt.
  \item Die Umkehrung (Teil~(ii)) erfordert BSD f\"ur $E_n$:
    Aus $A(n)=2B(n)$ folgt $L(E_n,1)=0$, und BSD impliziert
    $\mathrm{rang}(E_n) \ge 1$, also ist $n$ kongruent.
\end{enumerate}
Vollst\"andiger Beweis in \cite{Tunnell1983} und \cite{Koblitz1984}.
\end{proof}

% ================================================================
\section{Verifikation f\"ur kleine Werte}
% ================================================================

\begin{theorem}[Kongruente Zahlen: $n = 5, 6, 7$]
\label{thm:small_examples}
Die ganzen Zahlen $5, 6, 7$ sind kongruent.
\end{theorem}

\begin{proof}
\textbf{$n = 6$:} Dreieck $(3, 4, 5)$, Fl\"ache $= 6$.
Der zugeh\"orige Punkt auf $E_6: y^2 = x^3 - 36x$ ist $(25/4, 35/8)$.
Probe: $(35/8)^2 = 1225/64$ und $(25/4)^3 - 36(25/4)
= 15625/64 - 14400/64 = 1225/64$. \checkmark
Dieser Punkt hat unendliche Ordnung (er liegt nicht in der Torsionsuntergruppe
$\{O,(0,0),(\pm 6, 0)\}$).

\textbf{$n = 5$:} Dreieck $(3/2, 20/3, 41/6)$, Fl\"ache $= 5$.
Probe: $(3/2)^2 + (20/3)^2 = 9/4 + 400/9 = (81+1600)/36 = 1681/36 = (41/6)^2$. \checkmark
Fl\"ache $= (1/2)(3/2)(20/3) = (1/2)(10) = 5$. \checkmark

\textbf{$n = 7$:} Dreieck $(35/12, 24/5, 337/60)$, Fl\"ache $= 7$.
Probe: $(35/12)^2 + (24/5)^2 = 1225/144 + 576/25 = (30625 + 82944)/3600
= 113569/3600 = (337/60)^2$. \checkmark
Fl\"ache $= (1/2)(35/12)(24/5) = (1/2)(840/60) = (1/2)(14) = 7$. \checkmark
\end{proof}

\begin{remark}[Status von $n = 1, 2, 3$]
\begin{itemize}
  \item Fermat bewies mittels unendlichem Abstieg, dass $n=1$ nicht kongruent
    ist (\"aquivalent: $y^2 = x^4 - 1$ hat keine rationalen Punkte mit $y\ne 0$).
  \item $n = 2, 3$: Tunnells Satz sagt unter BSD vorher, dass $A(2) \ne 2B(2)$
    und $A(3) \ne 2B(3)$, was $n=2,3$ nicht-kongruent bedeutet.
    Dies ist aber noch nicht bedingungslos bewiesen.
  \item Die vollst\"andige bedingungslose L\"osung des kongruenten
    Zahlenproblems f\"ur $n=1,2,3$ folgt aus einem Beweis von BSD
    f\"ur die entsprechenden Kurven $E_1, E_2, E_3$.
\end{itemize}
\end{remark}

% ================================================================
\section{Verbindung zu BSD}
% ================================================================

\begin{theorem}[BSD und das kongruente Zahlenproblem]
\label{thm:bsd_cong}
Das kongruente Zahlenproblem ist unter BSD vollst\"andig gel\"ost:
Eine quadratfreie positive ganze Zahl $n$ ist genau dann kongruent,
wenn Tunnells Kriterium gilt: $A(n)=2B(n)$ (ungerades $n$) oder
$C(n)=2D(n)$ (gerades $n$).
\end{theorem}

\begin{proof}
\textbf{Richtung 1} ($n$ kongruent $\Rightarrow$ Kriterium):
Dies ist Tunnells Satz~\ref{thm:tunnell}(i), bedingungslos.

\textbf{Richtung 2} (Kriterium $\Rightarrow$ $n$ kongruent), bedingt durch BSD:
Aus $A(n)=2B(n)$ folgt $L(E_n,1)=0$.
Unter BSD: $L(E_n,1)=0$ impliziert $\mathrm{rang}(E_n)\ge 1$.
Nach Satz~\ref{thm:equiv} ist $n$ damit kongruent.
\end{proof}

% ================================================================
\section{Schlussfolgerung}
% ================================================================

Das kongruente Zahlenproblem, eines der \"altesten in der Mathematik,
wird -- zumindest bedingt -- durch die Theorie elliptischer Kurven und
BSD elegant gel\"ost.

Die \"Aquivalenz zwischen $n$ kongruent und $\mathrm{rang}(E_n(\Q)) \ge 1$
(Satz~\ref{thm:equiv}) reduziert das Problem auf die Arithmetik der
Kurve $E_n: y^2 = x^3 - n^2 x$.
Tunnells Satz liefert dann ein explizites, berechenbares Kriterium
durch Repr\"asentationszahlen tern\"arer quadratischer Formen.

Ein bedingungsloser Beweis, dass $n = 1, 2, 3$ nicht kongruent sind,
folgt aus einem Beweis von BSD f\"ur diese Kurven.

% ================================================================
\begin{thebibliography}{10}

\bibitem{Tunnell1983}
J.~B.~Tunnell,
A classical Diophantine problem and modular forms of weight $3/2$,
\emph{Invent.\ Math.}\ \textbf{72} (1983), 323--334.

\bibitem{Koblitz1984}
N.~Koblitz,
\emph{Introduction to Elliptic Curves and Modular Forms},
Graduate Texts in Mathematics~97,
Springer, New York, 1984.

\bibitem{Silverman1992}
J.~H.~Silverman und J.~Tate,
\emph{Rational Points on Elliptic Curves},
Undergraduate Texts in Mathematics,
Springer, New York, 1992.

\bibitem{Silverman1986}
J.~H.~Silverman,
\emph{The Arithmetic of Elliptic Curves},
Graduate Texts in Mathematics~106,
Springer, New York, 1986.

\bibitem{Cassels1991}
J.~W.~S.~Cassels,
\emph{Lectures on Elliptic Curves},
Cambridge University Press, 1991.

\end{thebibliography}

\end{document}
